\documentclass[11pt]{article}
\usepackage[T1]{fontenc}
\usepackage{amsmath}
\usepackage[margin=25mm]{geometry}
\title{Local writing, familiar workflow}
\author{LocalTeX demo}
\date{}

\begin{document}
\maketitle

\section{A browser workspace for local LaTeX}
Keep your sources on your own computer and compile with your
existing TeX installation. LocalTeX brings the source, project
files and PDF together in a browser, without an account or
cloud project upload.

\section{Write, preview, navigate}
Search your source, follow references, and use completion for
commands and citation keys. Independent views let you work on
several parts of the same file while sharing one draft.

The editor uses a native browser text field, making it a
practical home for grammar tools as well as LaTeX. Grammarly
is optional and has its own service and privacy requirements.

\subsection{Equations stay close to the explanation}
A driven mode can be written as
\begin{equation}
  \hat{H} = -\Delta \hat{a}^{\dagger}\hat{a}
  + \frac{U}{2}\hat{a}^{\dagger 2}\hat{a}^{2}
  + F\bigl(\hat{a}^{\dagger}+\hat{a}\bigr).
  \label{eq:mode}
\end{equation}
Select a source location and use the sync arrow to find its
PDF position. Double-click PDF text to return to its source.

\section{Local files, flexible projects}
Open an existing folder, import a project ZIP, or create files
from the tree. Subfolders can contain standalone figures or
other TeX documents with their own PDF outputs.


\section{A workflow you control}
Compile the current document with Ctrl+Enter. Use
Ctrl+Shift+Enter for a separate build target when editing
an included chapter. Inspect compiler diagnostics or search
the raw log when something needs attention.

Save frequently. Disk history and browser draft recovery
provide extra protection, but regular backups still matter.
\end{document}
